%% presentation.tex
%%
%% Presentation of the course ``Developers and their motivations'' of 
%% the Official Master on Libre Software (URJC)
%% http://master.libresoft.es
%%

%%---------------------------------------------------------------------
%%---------------------------------------------------------------------

%
% Compilar a .pdf con LaTeX (pdflatex)
% Es necesario instalar Beamer (paquete latex-beamer en Debian)
%

%
% Gráficos:
% Los gráficos pueden suministrarse en PNG, JPG, TIF, PDF, MPS
% Los EPS deben convertirse a PDF (usar epstopdf)
%

\documentclass{beamer}
\usetheme{Warsaw}
\usebackgroundtemplate{\includegraphics[width=\paperwidth]{format/libresoft-bg.png}}
\usepackage[spanish]{babel}
\usepackage[utf8]{inputenc}
\usepackage{graphics}
\usepackage{amssymb} % Simbolos matematicos

%\definecolor{libresoftgreen}{RGB}{162,190,43}
%\definecolor{libresoftblue}{RGB}{0,98,143}

%\setbeamercolor{titlelike}{bg=libresoftgreen}

%% Metadatos del PDF.
\hypersetup{
  pdftitle={Introduction to developers and their motivations},
  pdfauthor={Felipe Ortega},
  pdfcreator={GSyC/Libresoft},
  pdfproducer=PDFLaTeX,
  pdfsubject={Introduction to MSWL_DM},
}
%%

\begin{document}

\title{Introduction to developers and their motivations}
% \subtitle{Presentation about Eric S. Raymond's paper}
\institute{\{felipe.ortega\}@urjc.es\\
DEIO}
\author{Felipe Ortega}
\date{October 11, 2012}

\frame{
\maketitle
\begin{center}
\includegraphics[width=6cm]{format/gsyc-urjc}
\end{center}
}


% Si el titulo o el autor se quieren acortar para los pies de página
% se pueden redefinir aquí:
%\title{Titulo corto}
%\author{Autores abreviado}


%% LICENCIA DE REDISTRIBUCION DE LAS TRANSPAS
\frame{
~
\vspace{4cm}

\begin{flushright}
{\tiny
(cc) 2012 Felipe Ortega\\
  Some rights reserved. This work licensed under Creative Commons
  Attribution-ShareAlike License. To view a copy of full license, see
  http://creativecommons.org/licenses/by-sa/3.0/ or write to
  Creative Commons, 559 Nathan Abbott Way, Stanford,
  California 94305, USA.
%  Este documento (o uno muy similar) está disponible en \\
%  \url{http://gsyc.escet.urjc.es/~jjamor/}
}
\end{flushright}
}
%%

\section{Presentation of the course}

%%---------------------------------------------------------------

\begin{frame}
\frametitle{Administrative data}

\begin{itemize}
\item First semester, 3 ECTS credits
\item Teachers:
  \begin{itemize}
  \item Felipe Ortega (felipe.ortega at urjc.es)
  \item Departamento de Estadística e Investigación Operativa
  \item Office B-010 (Library) (Fuenlabrada campus)
  \item Daniel Izquierdo (dizquierdo at libresoft.es)
%   \item Juan José Amor (jjamor at libresoft.es)
  \item Departamento de Sistemas Telemáticos y Computación (GSyC)
  \item Office B-102 Biblioteca (Library) (Fuenlabrada campus)
  \end{itemize}
\item Schedule: Thursday, 17:00-21:00
\item Sessions: October 11, October 18, October 25, November 8 and November 15.
  \begin{itemize}
  \item Classroom 227, Departamental I, Fuenlabrada campus
%   \item Lab 7, Laboratorios II, Fuenlabrada campus
  \end{itemize}
\item Moodle course (\textbf{please, join it as soon as possible}): \\
  \url{http://gsyc.escet.urjc.es/moodle/course/view.php?id=123}
\end{itemize}
\end{frame}

%%---------------------------------------------------------------

\begin{frame}
\frametitle{Main topics in this course}

\begin{itemize}
\item Typical profiles of libre software developers
\item Motivations to participiate in libre software projects
\item Roles of participants and leadership in libre software projects
\item Temporal evolution of participation in libre software projects
\item Integration programs in libre software projects
\end{itemize}

\begin{quote}
\begin{center}
Understand who contributes to libre software development and why
\end{center}
\end{quote}

\end{frame}

%%---------------------------------------------------------------

\begin{frame}
\frametitle{Evaluation}

Please, refer to the course syllabus for specific grading info:
\begin{itemize}
\item Exercises (in forum)
\item Blog entries
\item Collaborative notebook
\item Specific report
\item ...
\end{itemize}

Deadline for submitting your work: \textbf{November 30th, 2012}.

\end{frame}

%%---------------------------------------------------------------

\begin{frame}
\frametitle{Final work: specific report}

\begin{itemize}
\item Report about FLOSS leader.
\item About 8-10 pages long.
\item Should cover personal and professional profile.
\item If involved in multiple projects, choose one or offer global retrospective.
\item \textit{Optional}: find a video of a talk or remarkable intervention (blog post, messages in mailing lists) of your target leader 
and comment on it.
\end{itemize}

\end{frame}

%%---------------------------------------------------------------

% \begin{frame}
% \frametitle{}
% 
% \begin{itemize}
% \item 
% \end{itemize}
% 
% \end{frame}

%%---------------------------------------------------------------

\begin{frame}
\frametitle{Some references}

\begin{itemize}
\item Web site of this course \\
  \url{http://gsyc.escet.urjc.es/moodle/course/view.php?id=123}
\item The cathedral and the bazaar (book of essays) \\
  \url{http://www.catb.org/esr/writings/cathedral-bazaar/}
\item Open Sources 2.0 The Continuing Evolution \\
  \url{http://oreilly.com/catalog/9780596008024/}
\end{itemize}

\end{frame}

\end{document}
